\section{Routes That Are Now Classified}
\label{sec:routes-classified}

Complete-period counting, native-period Bessel bounds, and separate capacity
envelopes do not resolve the late short-window problem. Accepted-anchor
recursion also returns the existing summatory coprime discrepancy after exact
CRT cancellation. These facts do not refute the quadratic survival condition;
they identify which additional information it must use.

The live twin-prime program is now narrow:
\begin{equation*}
\begin{aligned}
& \text{control }V_r\text{ relatively in the actual short window,}\\
& \text{control partial blocks,}
\quad
\text{then compose the signed layers.}
\end{aligned}
\end{equation*}

More optimization of unsigned capacity or complete-period norms cannot supply
those missing facts.

\subsection{The Fixed-Seed Scale Conflict}
\label{sec:fixed-seed-scale-conflict}

The exhaustion argument of Section~6 shows that no unsigned-capacity envelope can
certify survival. There is a deeper, structural reason complete-period
counting cannot place a survivor in a chosen window: the primorial modulus
outgrows the window itself.

Constrain the consecutive prime chain to end below the square horizon of its
seed prime $p$, so that the future head $Q$ satisfies $Q<p^2$. This keeps
the chain numerically short, but it conflicts with using many repeated
copies of one fixed seed residue inside the final safe window $[Q,Q^2)$.
The seed period is the primorial
\begin{equation*}
M_p=\prod_{r<p}r,
\end{equation*}

and by the prime number theorem in Chebyshev-theta form,
\begin{equation*}
\log M_p=\sum_{r<p}\log r\sim p.
\end{equation*}

Hence $M_p=\exp((1+o(1))p)$. Meanwhile $Q<p^2$ implies $Q^2<p^4$. Therefore
\begin{equation*}
\begin{aligned}
\frac{M_p}{Q^2}
&>\frac{\exp((1+o(1))p)}{p^4}
\longrightarrow\infty,
&&[\text{Asymptotic Limit}]\\
M_p&>Q^2
&&[\text{Eventually}].
\end{aligned}
\end{equation*}

For all sufficiently large scenarios satisfying $Q<p^2$, one fixed residue
class modulo $M_p$ occurs at most once in $[Q,Q^2)$. Its exact global
repetition frequency therefore cannot force local placement.

This is not a disproof of survival. It says that a proof cannot simultaneously
rely on a short chain under $p^2$ and on many local copies of one fixed seed.
A viable averaging argument must instead range over all seed residues at
once, use a much earlier seed with a longer filter chain, average over final
heads $Q$, or introduce an additional factorization or additive structure
that creates a bilinear variable. The prime number theorem is an explicit
external dependency in this argument.

This scale conflict is the geometric counterpart of the exhaustion verdict:
complete-period identities are exact globally but become locally uninformative
once the primorial exceeds the window. The signed local estimates of Sections~7--9
are precisely the response to this obstruction.
